\section{De la fonction image à l'observation numérique}
\label{sec:images}

\subsection{Domaine, intensités et échantillonnage}
On considère un domaine rectangulaire
$\Omega=(0,L_x)\times(0,L_y)\subset\R^2$ et une image propre
$\utrue:\Omega\to[0,1]$. Il s'agit d'une intensité scalaire en niveaux
de gris, non d'une image couleur. La normalisation fixe l'unité de
mesure ; elle ne rend pas le phénomène physique de bruit universel.
La représentation continue facilite l'écriture des opérateurs et des
conditions aux limites. Une image comportant des sauts n'est toutefois
pas une fonction partout différentiable.

Dans le calcul, l'image est une matrice de $N_y\times N_x$ pixels. Les
valeurs sont interprétées aux centres d'une grille uniforme, comme
précisé en section~\ref{sec:schemes}. Avec les pas unitaires utilisés,
$L_x=N_x$ et $L_y=N_y$ dans la convention pixel. Les coordonnées
normalisées employées pour construire les motifs synthétiques sont
une recette de génération ; elles ne remplacent pas les pas spatiaux
du solveur par $1/127$.
Les valeurs de référence sont construites dans $[0,1]$ avant tout
ajout de bruit. Cette propriété ne signifie pas que chaque image utilise
les deux extrémités de l'intervalle, ni qu'une méthode doit forcer ses
sorties à y rester. Les mêmes tableaux propres sont utilisés pour
toutes les réalisations d'un cas.

\subsection{Bruit gaussien : une hypothèse sur les pixels}
L'observation est écrite
\begin{equation}
 f=\utrue+\eta,\qquad
 f_{ij}=(\utrue)_{ij}+\sigma Z_{ij},\qquad
 Z_{ij}\overset{\mathrm{iid}}{\sim}\mathcal N(0,1).
 \label{eq:noise-model}
\end{equation}
L'indépendance concerne les pixels au sein d'un champ tiré. La graine
détermine une réalisation reproductible du générateur NumPy. Les
couplages entre cas dus à la réutilisation des graines sont décrits en
section~\ref{sec:protocol}. L'écriture fonctionnelle $f=\utrue+\eta$
est une notation de modélisation : le code produit un bruit fini sur
la grille. Il ne construit pas une fonction continue de bruit blanc,
dont la définition mathématique nécessiterait un autre cadre.

Même si chaque $\eta_{ij}$ a une espérance nulle, la moyenne du champ
réalisé n'est pas exactement nulle. Sa conservation par les solveurs
ne garantit donc pas une restauration exacte de la moyenne propre.
Le paramètre $\sigma$ est l'écart-type d'intensité par pixel. Il vaut
$0{,}025$, $0{,}05$ ou $0{,}10$ dans l'étude étendue.
Une variable gaussienne n'est pas bornée. Certaines observations
peuvent ainsi être négatives ou dépasser un. Les tronquer modifierait
la distribution, notamment près des niveaux sombres et clairs, puis
changerait les données et les métriques. Ni l'observation ni les états
des solveurs ne sont donc tronqués dans le calcul scientifique.
L'affichage des \emph{intensités} utilise seulement une échelle commune
$[0,1]$ ; les tableaux évalués et archivés restent inchangés. Les cartes
de gradient ont une échelle distincte explicitée dans leur légende.

\subsection{Donnée initiale, frontière et temps}
Les deux évolutions commencent au même état :
\begin{equation}
 u(x,y,0)=f(x,y).
 \label{eq:initial-data}
\end{equation}
Un flux nul est imposé sur la frontière. Pour la chaleur, il s'écrit
$\partial_{\boldsymbol n}u=0$ ; pour le modèle non linéaire,
$c(|\nabla u|)\nabla u\cdot\boldsymbol n=0$. Aucune image périodique
ni information extérieure n'est supposée. Sur la grille, la condition
est réalisée par l'absence exacte d'interfaces sortantes, et non par
l'ajout d'une bande de valeurs fixes.
Les temps archivés sont $t_n=n\dt$, avec $\dt=0{,}20$ et
$\dx=\dy=1$. Ils mesurent une échelle de diffusion dans cette convention
numérique, pas une durée en secondes. Le temps machine est une autre
variable, mesurée par chronomètre pour l'exécution informatique. Une
même trajectoire peut avoir le même temps de diffusion et des durées
informatiques différentes selon l'environnement.
